% figures/tid_vs_seu.tex — TID vs SEU 散点图（TID 悖论可视化）
% 数据来源：M2 radiation_dose.py

\begin{tikzpicture}
\begin{axis}[
    width=0.85\textwidth,
    height=9cm,
    xlabel={TID 中值 (\kradyr)},
    ylabel={SEU 相对风险 (log10, LEO=0)},
    xmin=0, xmax=140,
    ymin=-0.5, ymax=6.0,
    scatter/classes={
        leo={mark=square*,leoorange},
        meo={mark=triangle*,meocolor},
        geo={mark=diamond*,geocolor},
        heo={mark=pentagon*,heocolor},
        l1l2={mark=*,l1l2color},
        l4l5={mark=star,l4l5color}
    },
    grid=major,
    major grid style={dashed,gray!30},
    nodes near coords,
    every node near coord/.append style={font=\footnotesize, anchor=west},
]

% LEO 基准
\addplot[scatter, only marks, scatter src=explicit symbolic]
    coordinates {
        (20.0, 0.0) [leo]
    };
\node[anchor=south west, font=\footnotesize] at (axis cs:20.0,0.0) {LEO (基准)};

% MEO - SEE主导
\addplot[scatter, only marks, scatter src=explicit symbolic]
    coordinates {
        (10.0, 5.00) [meo]
    };
\node[anchor=west, font=\footnotesize] at (axis cs:10.0,5.00) {MEO};

% GEO - 混合型
\addplot[scatter, only marks, scatter src=explicit symbolic]
    coordinates {
        (6.5, 0.74) [geo]
    };
\node[anchor=south west, font=\footnotesize] at (axis cs:6.5,0.74) {GEO};

% HEO - 双杀
\addplot[scatter, only marks, scatter src=explicit symbolic]
    coordinates {
        (125.0, 5.50) [heo]
    };
\node[anchor=south east, font=\footnotesize] at (axis cs:125.0,5.50) {HEO};

% 地日L1/L2 - SEE主导 (低TID)
\addplot[scatter, only marks, scatter src=explicit symbolic]
    coordinates {
        (3.5, 2.50) [l1l2]
    };
\node[anchor=south east, font=\footnotesize, text=l1l2color] at (axis cs:3.5,2.50) {地日L1/L2};

% 地月L4/L5
\addplot[scatter, only marks, scatter src=explicit symbolic]
    coordinates {
        (10.0, 1.35) [l4l5]
    };
\node[anchor=east, font=\footnotesize, text=l4l5color] at (axis cs:10.0,1.35) {地月L4/L5};

\end{axis}
\end{tikzpicture}
